% ============================================================================
\chapter{Node.js : serveur, CLI, scripts d'automatisation}
\label{ch:node}
% ============================================================================

\begin{objectif}
\begin{itemize}
  \item Utiliser les modules de la biblioth\`eque standard~: \texttt{fs}, \texttt{path}, \texttt{os},
        \texttt{crypto}, \texttt{child\_process}, \texttt{worker\_threads}.
  \item \'Ecrire un serveur HTTP nu, puis un serveur Express complet~: routes, middleware,
        validation, erreurs, logs, upload.
  \item Structurer une application Node en couches, avec configuration par environnement.
  \item \'Ecrire des scripts CLI utiles (et les rendre robustes~: signaux, exit codes, stdout/stderr).
\end{itemize}
\end{objectif}

\section{Le mod\`ele, en trois phrases}

Node ex\'ecute du JavaScript sur V8 avec une boucle d'\'ev\'enements non bloquante~: les entr\'ees/sorties
sont d\'el\'egu\'ees \`a une pool de threads (libuv) et leurs rappels sont rejou\'es sur le fil principal.
Un processus ne g\`ere donc \emph{qu'un calcul \`a la fois}~: pour utiliser plusieurs coeurs, on lance
plusieurs processus (\verb|cluster|, \verb|pm2 -i max|, ou un conteneur par instance derri\`ere nginx).

\section{La biblioth\`eque standard qui compte}

\begin{codejs}
// ---- path, fs, os
import { join, dirname, basename, extname, resolve } from "node:path";
import { readFile, writeFile, readdir, mkdir, stat, rm, cp, appendFile, access } from "node:fs/promises";
import { existsSync, createReadStream, createWriteStream, watch } from "node:fs";
import { homedir, tmpdir, cpus, totalmem } from "node:os";

const here = dirname(new URL(import.meta.url).pathname);
const data = await readFile(join(here, "data.json"), "utf8");
await writeFile(join(here, "out.txt"), "ok\n", { flag: "a" });
await mkdir(join(here, "a/b/c"), { recursive: true });
await cp("src", "backup", { recursive: true, force: true });
for (const entree of await readdir(here, { withFileTypes: true })) {
  if (entree.isDirectory() && !entree.name.startsWith(".")) console.log(entree.name);
}
await access("fichier.txt").then(() => console.log("lisible")).catch(() => console.log("absent ou interdit"));

// GROS fichier ligne a ligne sans le charger en memoire :
import { createInterface } from "node:readline";
const rl = createInterface({ input: createReadStream("gros.log"), crlfDelay: Infinity });
let n = 0;
for await (const ligne of rl) if (/ERR/.test(ligne)) n++;
console.log(n, "lignes d'erreur");

// GROS ecriture en flux (les backpressures sont gerees par pipeline)
import { pipeline } from "node:stream/promises";
import { createGzip } from "node:zlib";
await pipeline(createReadStream("gros.csv"), createGzip(), createWriteStream("gros.csv.gz"));

// Surveillance de fichiers
const w = watch(here, { persistent: false }, (evt, filename) => console.log(evt, filename));
w.close();

// ---- crypto : hachage, jetons, chiffrement
import { createHash, randomBytes, createHmac, randomUUID, scrypt, timingSafeEqual } from "node:crypto";
createHash("sha256").update("mot de passe + sel").digest("hex");
randomUUID();                       // "d94b..." (v4)
randomBytes(32).toString("base64url");   // jeton d'API
const sel = randomBytes(16);
const cle = await promisify(scrypt)(mdp, sel, 64);   // derive de cle (en prod : argon2 ou bcrypt)
timingSafeEqual(a, b);              // comparaison sans fuite temporelle

// ---- child_process
import { execFile, spawn } from "node:child_process";
import { promisify } from "node:util";
const run = promisify(execFile);
const { stdout, stderr } = await run("git", ["rev-parse", "--short", "HEAD"]);
const ff = spawn("ffmpeg", ["-i", "in.mov", "out.mp4"], { stdio: "inherit" });
ff.on("exit", (code, sig) => console.log("fini", code, sig));

// ---- process : signaux, code de sortie, arguments
const args = process.argv.slice(2);
const opts = Object.fromEntries(args.filter((a) => a.startsWith("--")).map((a) => {
  const [k, v = "true"] = a.replace(/^--/, "").split("=");
  return [k, v === "true" ? true : v];
}));
process.on("SIGINT", () => arretPropre(130));
process.on("SIGTERM", () => arretPropre(143));
async function arretPropre(code) {
  await Promise.allSettled([serveur.closeAsync?.(), db.end()]);
  process.exitCode = code;      // ne pas appeler process.exit brutalement : on perd les bufferisations
}
console.error("aller sur stderr"); process.exitCode = 1;   // exit code correct pour les scripts
\end{codejs}

\section{Le serveur HTTP nu, pour comprendre}

\begin{codejs}
import { createServer } from "node:http";

const serveur = createServer(async (req, res) => {
  const url = new URL(req.url, `http://${req.headers.host}`);
  const json = (status, objet) => {
    const corps = JSON.stringify(objet);
    res.writeHead(status, {
      "content-type": "application/json; charset=utf-8",
      "content-length": Buffer.byteLength(corps),
      "x-powered-by": "node nu",
    });
    res.end(corps);
  };
  try {
    if (url.pathname === "/api/sante") return json(200, { ok: true, uptime: process.uptime() });
    if (url.pathname === "/api/echo" && req.method === "POST") {
      const chunks = [];
      for await (const c of req) chunks.push(c);
      return json(200, { recu: JSON.parse(Buffer.concat(chunks).toString() || "{}") });
    }
    json(404, { erreur: `route inconnue ${req.method} ${url.pathname}` });
  } catch (e) {
    json(400, { erreur: e.message });
  }
});
serveur.listen(3000, () => console.log("http://localhost:3000"));
serveur.keepAliveTimeout = 65_000;            // > au timeout du reverse proxy (sinon 502 aleatoires)
serveur.requestTimeout = 300_000;
process.on("SIGTERM", () => serveur.close(() => process.exit(0)));
\end{codejs}

\section{Express~: le cadre minimaliste de r\'ef\'erence}

Documentation~: \link{expressjs.com/fr}{https://expressjs.com/fr/}. Version~5 (2024+) g\`ere
enfin les erreurs des fonctions \texttt{async} sans encapsulation.

\begin{terminal}
mkdir api && cd api && npm init -y
npm pkg set type=module
npm i express helmet cors compression morgan zod dotenv
npm i -D nodemon eslint
\end{terminal}

\begin{codejs}
// src/app.js
import express from "express";
import helmet from "helmet";
import cors from "cors";
import morgan from "morgan";
import { routers } from "./routes/index.js";
import { environement } from "./config.js";
import { gestionErreur } from "./erreurs.js";

export function creerApp() {
  const app = express();
  app.disable("x-powered-by");
  app.set("trust proxy", 1);                       // derriere nginx : lire X-Forwarded-*
  app.use(helmet({ crossOriginResourcePolicy: { policy: "cross-origin" } }));
  app.use(cors({ origin: environement.origines, credentials: true }));
  app.use(express.json({ limit: "1mb" }));
  app.use(express.urlencoded({ extended: true, limit: "100kb" }));
  app.use(morgan(environement.prod ? "combined" : "dev", {
    skip: (req) => req.path === "/favicon.ico",
  }));
  app.use("/static", express.static("public", { maxAge: "7d", immutable: true }));

  routers.forEach(([chemin, r]) => app.use(chemin, r));
  app.use((req, res) => res.status(404).json({ erreur: "introuvable", chemin: req.path }));
  app.use(gestionErreur);
  return app;
}

// src/routes/articles.js
import { Router } from "express";
import { z } from "zod";
import { asyncHandler, HttpError } from "../erreurs.js";
import { depot } from "../depot.js";

const schemaArticle = z.object({
  titre: z.string().trim().min(3).max(160),
  contenu: z.string().trim().min(10),
  tags: z.array(z.string().min(1).max(30)).max(12).default([]),
  publie: z.boolean().default(false),
});

export const routeArticles = Router();

routeArticles.get("/", asyncHandler(async (req, res) => {
  const q = z.object({
    page: z.coerce.number().int().min(1).default(1),
    taille: z.coerce.number().int().min(1).max(100).default(20),
    recherche: z.string().trim().max(120).optional(),
  }).parse(req.query);
  const { items, total } = await depot.liste({ ...q });
  res.links({ next: total > q.page * q.taille ? `?page=${q.page + 1}` : undefined });
  res.json({ items, total, page: q.page });
}));

routeArticles.get("/:id(\\d+)", asyncHandler(async (req, res) => {
  const article = await depot.parId(Number(req.params.id));
  if (!article) throw new HttpError(404, "Article introuvable");
  res.set("etag", `W/"${article.maj}"`).vary("accept-encoding");
  res.status(200).json(article);
}));

routeArticles.post("/", asyncHandler(async (req, res) => {
  const donnees = schemaArticle.parse(req.body);
  const cree = await depot.creer(donnees);
  res.location(`/api/articles/${cree.id}`).status(201).json(cree);
}));

routeArticles.delete("/:id(\\d+)", asyncHandler(async (req, res) => {
  await depot.supprimer(Number(req.params.id));
  res.status(204).end();
}));

// src/erreurs.js
export class HttpError extends Error {
  constructor(status, message, details) { super(message); this.status = status; this.details = details; }
}
export const asyncHandler = (fn) => (req, res, next) => Promise.resolve(fn(req, res, next)).catch(next);
export function gestionErreur(err, req, res, _next) {
  if (err?.name === "ZodError") {
    return res.status(422).json({ erreur: "Donnees invalides", details: err.issues.map((i) => ({ champ: i.path.join("."), message: i.message })) });
  }
  const status = err.status ?? err.statusCode ?? 500;
  if (status >= 500) console.error({ msg: err.message, stack: err.stack, url: req.originalUrl, ip: req.ip });
  res.status(status).json({ erreur: status >= 500 ? "Erreur interne" : err.message });
}
\end{codejs}

\subsection{Configuration par environnement~: la bonne habitude d\`es le premier jour}
\label{sec:node-config}

\begin{codejs}
// src/config.js --- une seule porte d'entree, validee au demarrage
import "dotenv/config";
import { z } from "zod";

const schema = z.object({
  NODE_ENV: z.enum(["development", "test", "production"]).default("development"),
  PORT: z.coerce.number().int().default(3000),
  DATABASE_URL: z.string().url(),
  JWT_SECRET: z.string().min(32),
  LOG_LEVEL: z.enum(["debug", "info", "warn", "error"]).default("info"),
  ORIGINES: z.string().default("http://localhost:5173"),
  UPLOAD_MAX_MO: z.coerce.number().default(5),
});
const parsed = schema.safeParse(process.env);
if (!parsed.success) {
  console.error("Configuration invalide :", parsed.error.flatten().fieldErrors);
  process.exit(1);
}
export const environement = {
  ...parsed.data,
  prod: parsed.data.NODE_ENV === "production",
  origines: parsed.data.ORIGINES.split(",").map((s) => s.trim()),
};
\end{codejs}

Jamais de secret dans le code, dans Git, ni dans une variable expos\'ee au navigateur.
\verb|.env| pour le dev local, variables d'environnement r\'eelles en production
(\link{Doppler}{https://doppler.com/}, \link{1Password~CLI}{https://developer.1password.com/docs/cli/},
ou simplement le fichier d'environnement de systemd / de ton h\'ebergeur).

\section{Ce qu'Express ne fait pas pour toi (et qu'il faut ajouter)}

\begin{itemize}
  \item \textbf{Cookies et sessions}~: \link{cookie-parser}{https://github.com/pillarjs/cookie-parser},
        \link{express-session}{https://github.com/expressjs/session} (ou des sessions sign\'ees sans
        \'etat, cf. \cref{ch:js-secu}).
  \item \textbf{Upload}~: \link{multer}{https://github.com/expressjs/multer}, avec taille maximale,
        liste blanche de types et stockage hors de la racine web.
  \item \textbf{Limitation de d\'ebit}~: \link{express-rate-limit}{https://github.com/nickdesclos/express-rate-limit}.
  \item \textbf{Corr\'elation et logs structur\'es}~: \link{pino}{https://github.com/pinojs/pino}
        (JSON, tr\`es rapide) plut\^ot que \verb|console.log|.
  \item \textbf{T\^aches diff\'er\'ees}~: \link{BullMQ}{https://docs.bullmq.io/} sur Redis.
  \item \textbf{Temps r\'eel}~: \link{Socket.IO}{https://socket.io/fr/} ou WebSocket brut.
\end{itemize}

\begin{codejs}
// Exemple d'upload durci
import multer from "multer";
import { extname } from "node:path";
import { randomUUID } from "node:crypto";

const TYPES = new Map([["image/png", ".png"], ["image/jpeg", ".jpg"], ["image/webp", ".webp"]]);
const upload = multer({
  storage: multer.diskStorage({
    destination: "var/uploads",
    filename: (req, file, cb) => cb(null, `${randomUUID()}${TYPES.get(file.mimetype) ?? ".bin"}`),
  }),
  limits: { fileSize: 5 * 1024 * 1024, files: 1 },
  fileFilter: (req, file, cb) => cb(TYPES.has(file.mimetype) ? null : new Error("type refuse"), TYPES.has(file.mimetype)),
});
route.post("/avatar", upload.single("fichier"), asyncHandler(async (req, res) => {
  if (!req.file) throw new HttpError(422, "Fichier manquant ou type non autorise");
  res.json({ url: `/u/${req.file.filename}` });
}));
\end{codejs}

\section{Alternatives \`a Express, et quand les choisir}

\begin{tabz}{lXl}
  \headrow \thd{Cadre} \& \thd{Id\'ee directrice} \& \thd{Lien} \\
  \midrule
  \link{Hono}{https://hono.dev/} \& ultra-l\'eger, tourne sur Node, Bun, Deno, Cloudflare Workers, navigateur. Mon choix pour une API neuve. \& 14 kB \\
  \link{Fastify}{https://fastify.dev/} \& plus rapide qu'Express, sch\'emas JSON int\'egr\'es, plugins. \& $\approx$2x \\
  \link{Koa}{https://koajs.com/} \& le \og ~Express repens\'e~ \fg{} par les auteurs d'Express. \& minimaliste \\
  \link{NestJS}{https://nestjs.com/} \& architecture d\'ecorateurs/injection, tr\`es TypeScript, \'equipes Java qui passent au Node. \& complet \\
  \link{Next.js}{https://nextjs.org/} \& React fullstack~: pages, API routes, rendu serveur. \& cadre UI \\
  \link{Express~5}{https://expressjs.com/fr/} \& le standard historique, documentation immense. \& classique \\
\end{tabz}

\section{Un CLI utile en 60 lignes}

\begin{codejs}
#!/usr/bin/env node
// scripts/rename.js --- renomme en lot : node rename.js --dossier ./photos --prefix 2026 --debut 1
import { readdir, rename, stat } from "node:fs/promises";
import { join, extname } from "node:path";

const opts = Object.fromEntries(process.argv.slice(2).map((a) => {
  const [k, v = "true"] = a.replace(/^--/, "").split("=");
  return [k, v];
}));

if (opts.aide) {
  console.log(`Usage : rename --dossier CHEMIN --prefix TEXTE [--debut N] [--sec]

  --dossier  dossier a traiter (obligatoire)
  --prefix   debut du nouveau nom
  --debut    numero de depart (defaut 1)
  --sec      applique vraiment (sinon : simulation)`);
  process.exit(0);
}
const dossier = opts.dossier ?? process.exit(fail("--dossier est requis"));
const prefix = opts.prefix ?? "fichier";
let compteur = Number(opts.debut ?? 1);

const fichiers = (await readdir(dossier))
  .filter((f) => !f.startsWith("."))
  .sort((a, b) => a.localeCompare(b, "fr", { numeric: true }));

for (const f of fichiers) {
  const neuf = `${prefix}-${String(compteur++).padStart(3, "0")}${extname(f).toLowerCase()}`;
  if (neuf === f) continue;
  const { size, mtime } = await stat(join(dossier, f));
  console.log(`${f}  ->  ${neuf}   (${(size / 1024).toFixed(0)} Ko, ${mtime.toISOString().slice(0, 10)})`);
  if (opts.sec) await rename(join(dossier, f), join(dossier, neuf));
}
if (!opts.sec) console.log("\nSimulation. Relance avec --sec pour appliquer.");

function fail(msg) { console.error(`Erreur : ${msg}`); return 1; }
\end{codejs}

\subsection{Automatiser~: Makefile et scripts de d\'eveloppement}

\begin{codeconf}
# Makefile --- la interface unique du projet (fais `make` pour voir la liste)
.PHONY: help dev check test lint db-reset docker-up docker-down

help:            ## affiche cette aide
	@grep -E '^[a-zA-Z_-]+:.*?## ' $(MAKEFILE_LIST) | awk 'BEGIN {FS = ":.*?## "}; {printf "  \033[36m%-12s\033[0m %s\n", $$1, $$2}'
dev:             ## lance API + site en parallle
	@npm run dev
check: lint test ## tout verifier (lint + typecheck + tests)
	npm run typecheck
lint:            ## analyse statique
	npm run lint -- --fix
test:            ## tests unitaires
	npm test
db-reset:        ## recree la base et la remplit
	node scripts/migrate.js --fresh && node scripts/seed.js
docker-up:       ## demarre les services (MySQL, Redis, Mailpit)
	docker compose up -d
docker-down:     ## arrete tout
	docker compose down
\end{codeconf}

\begin{codejson}
{
  "name": "projet-node",
  "scripts": { "dev": "node --watch --env-file=.env src/server.js" },
  "devDependencies": { "nodemon": "^3.1.0" }
}
\end{codejson}

Depuis Node 18.11, \cmd{node --watch} suffit au quotidien. Pour recharger aussi au changement de
fichier non-JS (views, .env), \link{nodemon}{https://nodemon.io/} reste pratique. Pour red\'emarrer
proprement un processus en production~: \link{pm2}{https://pm2.keymetrics.io/} (ou mieux~: systemd).

\begin{terminal}
# pm2 en production (si tu n'as pas de orchestrateur)
npm i -g pm2
pm2 start src/server.js -i max --name api --node-args="--max-old-space-size=1024"
pm2 logs api --lines 100
pm2 monit ; pm2 reload api ; pm2 save ; pm2 startup
\end{terminal}

\begin{exercice}
\begin{enumerate}
  \item \'Ecris un serveur HTTP nu qui sert un fichier statique depuis \verb|public/| avec les
        en-t\^etes corrects~: \verb|content-type| (par extension), \verb|cache-control|, ETag
        bas\'e sur taille+date, et gestion du \verb|If-None-Match| (304).
  \item Reprends-le en Express, en ajoutant~: route \verb|/api/sante|, validation par
        \verb|zod| d'un POST, gestion d'erreur centralis\'ee, et logs \verb|pino|.
  \item Construis une API CRUD compl\`ete pour un blog (articles + commentaires) stock\'ee sur
        disque en JSON, avec pagination, recherche, et code de sortie correct partout.
        \emph{Ce sera ta base de comparaison avec les versions PHP et Java.}
  \item \'Ecris un CLI \verb|depense|~: \verb|add "cafe" 2.5 --cat=alimentation|,
        \verb|list --month=2026-01|, \verb|total --by=category|, donn\'ees dans
        \verb|~/.depenses.json|. Ajoute \verb|--json| pour la sortie machine.
  \item \'Ecris un script qui v\'erifie la sant\'e de ton API (5 requ\^etes, timeouts, format) et
        retourne un code de sortie non nul en cas d'\'echec~: utilisable dans un cron ou une CI.
  \item Lance deux instances derri\`ere \cmd{node --cluster} (ou l'option \verb|-i max| de pm2)
        et mesure le d\'ebit avec \link{autocannon}{https://github.com/mcollina/autocannon} :
        explique le r\'esultat avec le mod\`ele mono-fil.
\end{enumerate}
\end{exercice}

\begin{ressources}
\begin{itemize}
  \item Docs Node (la r\'ef\'erence absolue) --- \url{https://nodejs.org/fr/docs/latest/api/}
  \item Node --- \emph{Learn JavaScript on the server} --- \url{https://nodejs.org/learn}
  \item Express guide --- \url{https://expressjs.com/fr/guide/routing.html}
  \item \link{Node Best practices (Goldberd)}{https://github.com/goldberio/express-nodejs-best-practices}
  \item \link{Node Security Workshop}{https://github.com/ThankyouForComing/Node-Security-Workshop} --- et \url{https://snyk.io/fr/blog/}
\end{itemize}
\end{ressources}

\index{Node.js}\index{Express}\index{CLI}\index{serveur}
